%-*-latex-*-
%-*-latex-*-
\newcommand\COURSE{ciss245}
\newcommand\ASSESSMENT{q2903}
\newcommand\ASSESSMENTTYPE{Quiz}
\newcommand\POINTS{\textwhite{xxx/xxx}}

\input{myquizpreamble}
\input{yliow}
\input{\COURSE}
%-*-latex-*-

%-*-latex-*-

%-*-latex-*-
\renewcommand\TITLE{\ASSESSMENTTYPE \ \ASSESSMENT}

\renewcommand\EMAIL{}
\renewcommand\AUTHOR{eagraham1@cougars.ccis.edu}

\textwidth=6in
\begin{document}
\topmatterthree

%------------------------------------------------------------------------------
\nextq
Complete the following to store value \verb!v! at the integers
from address \verb!start! up to but not including \verb!end!.
Recurse on \verb!start!, i.e., in your recursive case for
\verb!fill(start, end, v)!, you should call
\verb!fill(start + 1, end, v)!.
(5 seconds. Test.)
\\
\ANSWER
\begin{answercode}
void fill(int * start, * end, int v)
{
  if(start == end) return;

  *start = v;
  
  fill(start + 1, end, v);
}
\end{answercode}

%------------------------------------------------------------------------------
\nextq
This is the same as above, except that you recurse on \verb!end!,
i.e., in your recursive case for \verb!fill(start, end, v)!, you should call
\verb!fill(start, end - 1, v)!.
(5 seconds. Test.)
\\
\ANSWER
\begin{answercode}
void fill(int * start, * end, int v)
{
  if (start == end) return; 

  fill(start, end - 1, v);
  
  *(end - 1) = v; 
}
\end{answercode}

%------------------------------------------------------------------------------
\nextq
This is the same as above, except that you recurse using divide-and-conquer.
(20 seconds. Test.)
\\
\ANSWER
\begin{answercode}
void fill(int * start, * end, int v)
{
  if (start == end) return;
  
  int* mid = start + (end - start) / 2;
  
  fill(start, mid, v);
  fill(mid, end, v);

}
\end{answercode}

%------------------------------------------------------------------------------
\nextq
Complete the following function that counts the number of times \verb!target!
occurs in the integers
from address \verb!start! up to but not including \verb!end!.
Recurse at \verb!start!.
(5 seconds.)
\\
\ANSWER
\begin{answercode}
int count(int * start, int * end, int target)
{
  if(start == end) return 0;
  return (*start == target) + count(start + 1, end, target);
}
\end{answercode}

%------------------------------------------------------------------------------
\nextq
Complete the following function that counts the number of times \verb!target!
occurs in the integers
from address \verb!start! up to but not including \verb!end!.
Recurse at \verb!end!.
(5 seconds.)
\\
\ANSWER
\begin{answercode}
int count(int * start, int * end, int target)
{
  if(start == end) return 0;
  end--;

  return (*end == target) + count(start, end, target);  
}
\end{answercode}

%------------------------------------------------------------------------------
\nextq
Complete the following function that counts the number of times \verb!target!
occurs in the integers
from address \verb!start! up to but not including \verb!end!.
You must use divide-and-conquer.
(20 seconds.)
\\
\ANSWER
\begin{answercode}
int count(int * start, int * end, int target)
{
  if (start >= end) return 0;

  int * mid = start + (end - start) / 2;
  return count(start, mid, target) + count(mid, end, target);
}
\end{answercode}

%------------------------------------------------------------------------------
\nextq
Complete the \verb!bubblesort! function and the \verb!bubblesort_pass!
function.
This was already covered in class. So this is a freebie.
Finish the following bubblesort function by completing all the functions using
recursion.
Do not look at your class notes and try to do is on your own.
If you are lost, \textit{then} look at your notes.
There must be no loops in the program.
Your \verb!bubblesort! must call the \verb!bubblesort_pass! function that
executes one pass of the \verb!bubblesort!.
\\
\ANSWER
\begin{answercode}
#include <iostream>

void swap(int * x, int * y)
{
  int * temp = *x;
  *x = *y;
  *y = *temp;
}

void bubblesort_pass(int * start, int * end)
{
  if (start >= end - 1) return;

  if (*start > *(start + 1)) swap(start, start + 1);

  bubblesort_pass(start + 1, end);
}

void bubblesort(int * start, int * end)
{
  if (start >= end - 1) return;

  bubblesort_pass(start, end);
  bubblesort(start, end - 1);
}

void println(int * start, int * end)
{
    if (start >= end - 1)
    {
        std::cout << '\n';
    }
    else
    {
        std::cout << *start << ' ';
        println(start + 1, end);
    }
}

int main()
{
    int x[] = {9, 7, 5, 3, 1, 0, 2, 4, 6, 8};
    bubblesort(&x[0], &x[10]);
    println(&x[0], &x[10]);
    return 0;
}
\end{answercode}

%------------------------------------------------------------------------------
%-*-latex-*-
\newpage
\input{instructions}
\end{document}

